\documentclass[10pt,a4paper]{amsart}
\usepackage[margin=19mm]{geometry}
\usepackage{amsmath,amssymb,mathtools}
\usepackage[scr=rsfs,cal=euler]{mathalfa}
\usepackage{xfrac}
\usepackage[unicode,hidelinks,bookmarks=false]{hyperref}
\newcommand{\ie}{i.e.\ }
\newcommand{\cf}{cf.\ }
\newcommand{\resp}{respectively\ }
\newcommand{\Subsection}{\subsection}
\providecommand{\nobreakdash}{\nobreak\hbox{-}}
\setlength{\emergencystretch}{2em}
\multlinegap0pt
\usepackage{amssymb}
\usepackage{enumerate}
\usepackage{tikz}
\newtheorem{lem}{Lemma}
\newcommand{\CZ}{\mathcal{CZ}}
\newcommand{\W}{\mathcal{W}}
\title{Local centralizer rigidity for twisted Weyl chamber flows}
\author{Zhijing Wendy Wang}
\date{Source: arXiv:2606.09075v1 (CC BY 4.0)}
\begin{document}
\maketitle

\noindent\emph{Source and licence.} Local centralizer rigidity for twisted Weyl chamber flows. Version arXiv:2606.09075v1 (CC BY 4.0), distributed by arXiv under the Creative Commons Attribution 4.0 International licence, http://creativecommons.org/licenses/by/4.0/, as recorded in the arXiv licence metadata for this exact version and shown on the abstract page https://arxiv.org/abs/2606.09075v1. Full licence notice: \texttt{licenses/arxiv-2606.09075-CC-BY-4.0.txt}.\par\smallskip

\noindent\textit{Editorial scope.} Counted unit: the article's lem gcontr (source0.tex lines 1223--1232). The article's complete original proof of it is delivered with it (source0.tex lines 1234--1279, one \texttt{proof} environment of 46 lines). No mathematics is written, repaired or re-derived: the statement and the proof are the article's own lines, in the article's own notation, and no retained line is substituted or re-flowed.  No further source-local context is needed: every cross-reference of every retained line resolves inside the two delivered spans.  Nothing was omitted from any retained span, so the package carries no omission record.  No retained line contains a citation, so the wrapper carries no bibliography and no bibliography entry.  No retained line contains a figure environment, an image-inclusion command or drawing code, so no image is delivered and the package declares no asset dependency.  Editorial formatting only, in addition: an AMS \texttt{amsart} preamble that restates the article's own declarations and macros used by the retained lines, namely the environments \texttt{lem} on separate counters, together with the standard packages those lines need.  The retained environments print in this document as Lemma 1 in order of appearance, whereas the article prints the same items with the numbering of its own full document.  The complete native source of the article as distributed in the arXiv e-print for this exact version is bundled unchanged as \texttt{sources/202381/source0.tex}.
\begin{lem}\label{gcontr}
    Suppose $g \in \CZ_0(f)$ is projectively $\epsilon_0/2$-close to an element $g_0 \in \CZ_0(f_0) \simeq G^c$. For a root or weight $\mu\in \Lambda_0$ of the (twisted) Weyl chamber flow, suppose $\mathrm{sign}(\mu(g_0) + \nu_0(g_0)) = \mathrm{sign}(\mu(g_0))$ for any root or weight $\nu_0\in \Lambda_0$ with $\nu_0(f_0) = 0$, and suppose $g_0$ is not in the $\epsilon_0$-cone of the Weyl chamber wall $\ker(\mu)$. Then there exist $\epsilon > 0$, and $N > 0$, depending only on $g_0$, such that for every $x \in X$ and $y \in \W^\mu_\#(x)$ with $d(x,y) < \epsilon$, and for any $n > N$, we have:
    \[
    d(g^n(x), g^n(y)) \le e^{n \theta \mu(g_0)/4} d(x,y)^{\theta^2}, \quad \text{if } \mu(g_0) < 0,
    \]
    \[
    d(g^{-n}(x), g^{-n}(y)) \le e^{-n \theta \mu(g_0)/4} d(x,y)^{\theta^2}, \quad \text{if } \mu(g_0) > 0.
    \]
    Here, $\theta < 1$ is the bi-Hölder constant of the leaf conjugacy $\phi$.
\end{lem}

\begin{proof}[Proof of Lemma \ref{gcontr}]


    We restrict our discussion to the case $\mu(g_0) < 0$, where $g_0$ uniformly contracts $\W^\mu_{f_0}$. The case $\mu(g_0) > 0$ follows by considering $g^{-1}$.

By the definition of projective closeness, for every \(z\in X\), the
center displacement
\[
\tilde g(z)z^{-1}\in G^c
\]
has Cartan projection lying in the prescribed \(\epsilon_0/2\)-cone about
\(g_0\). Hence, for each \(j\ge0\), the one-step displacement
\[
\tilde g(\tilde g^j z)(\tilde g^j z)^{-1}
\]
satisfies the same cone condition. Applying the corresponding one-step
estimate successively along the orbit \(z,\tilde g z,\ldots,\tilde g^{n-1}z\),
we obtain for any $w \in \W^\mu_{f_0}(z)$,
\[
d(\tilde g^n z,\tilde g^n w)
   \le e^{n\mu(g_0)/2}d(z,w),
\qquad w\in \W^\mu_{f_0}(z).
\]

    Furthermore, since the leaves $\W^c_{f_0}$ and $\W^\mu_{f_0}$ are cosets of subgroups of $G$ with orthogonal Lie algebras, there exists a constant $\epsilon > 0$ such that for any $z \in X$ and $w \in \W^\mu_{f_0}(z)$ with $d(z,w) < \epsilon$, we have:
    \[
    d(z, \tilde{\W}^c_{f_0}(w)) \le 2d(z,w).
    \]
    Since $f$ is a $C^1$-perturbation of $f_0$ and $\W^\mu_\#$ is a subfoliation of $\W^s_f$ or $\W^u_f$, a similar property holds for $f$: for any $z \in X$ and $w \in \W^\mu_\#(z)$ with $d(z,w) < \epsilon$, we have $d(z, \tilde{\W}^c_f(w)) \le 2d(z,w)$ in the cover $G$.

    Therefore, using the Hölder continuity of the leaf conjugacy $\phi$, for every $x \in X$ and every $y \in \W^\mu_\#$ with $d(x,y) \le \epsilon$, define:
    \[
    y' = \phi\big(\W^\mu_{f_0}(\phi^{-1}(x))\big) \cap \W^c_f(y), \quad y'_n = \phi(\W^\mu_{f_0}(\phi^{-1}(g^n(x))) \cap \W^c_f(g^n(y)).
    \]
    Then we have:
    \[
    \begin{aligned}
    d(g^n(x), g^n(y)) &\le 2 d(g^n(x), y'_n) \\
    &\le 2C d(\hat{g}^n(\phi^{-1}(x)), \phi^{-1}(y'_n))^\theta \\
    &\le 2C \left(e^{n\mu(g_0)/2} d(\phi^{-1}(x), \phi^{-1}(y'))\right)^\theta \\
    &\le 4C e^{n \theta \mu(g_0)/2} d(\phi^{-1}(x), \phi^{-1}(y))^\theta \\
    &\le 4C^2 e^{n \theta \mu(g_0)/2} d(x,y)^{\theta^2}.
    \end{aligned}
    \]
    By choosing $N$ sufficiently large, we obtain the desired result.
\end{proof}

\end{document}
